\documentclass[10pt,journal,compsoc]{IEEEtran}

\usepackage{amsmath,amssymb,amsfonts}
\usepackage{algorithmic}
\usepackage{graphicx}
\usepackage{textcomp}
\usepackage{xcolor}
\usepackage{booktabs}
\usepackage{hyperref}
\usepackage{listings}

\lstset{
  language=bash,
  basicstyle=\ttfamily\footnotesize,
  breaklines=true,
  frame=single,
  columns=fullflexible
}

\begin{document}

\title{Thermodynamically Consistent Inter-Cell Boundary Metrics for Discrete Global Grid Systems in Planetary Biospheric Simulation}

\author{Chief~Systems~Architect~and~The~Web~of~Life~Consortium%
\thanks{Source code repository: \url{https://github.com/pascalranoroarijaona/WebOfLife}}}

\markboth{Web of Life Technical Report, Sprint 051,~Vol.~51, No.~1, 2025}%
{Architect \MakeLowercase{\textit{et al.}}: H3CellInterfaceMetrics for Discrete Global Grid Systems}

\IEEEtitleabstractindextext{%
\begin{abstract}
Discrete Global Grid Systems (DGGS), specifically hexagonal hierarchical partitions such as Uber H3, offer near-isometric discretizations of spherical manifolds for planetary-scale earth system modeling. However, numerical solutions of partial differential equations governing hydrological, atmospheric, and biogeochemical mass-energy transfers require precise metric definitions across inter-cell Voronoi interfaces. In this paper, we present the theoretical formulation, architectural typing, and empirical verification of the \texttt{H3CellInterfaceMetrics} contract developed for the Web of Life simulation platform. The metric system formalizes shared Voronoi edge lengths, geodesic centroid distances, local tangent contact normals, topographic gradients, and vertically stratified contact cross-sections. We establish analytical proofs and machine-precision validations demonstrating that all inter-cell transfer operators parameterized by this metric satisfy First Law conservation invariance ($\Phi_{i \to j} = -\Phi_{j \to i}$) and Second Law non-negative entropy generation ($\dot{S}_{\text{gen}} \ge 0$).
\end{abstract}

\begin{IEEEkeywords}
Discrete Global Grid System, H3 Hexagonal Grid, Finite Volume Method, Thermodynamic Invariance, Computational Ecology, Earth System Modeling.
\end{IEEEkeywords}}

\maketitle
\IEEEdisplaynontitleabstractindextext
\IEEEpeerreviewmaketitle

\section{Introduction}
\IEEEPARstart{P}{lanetary-scale} biospheric and biogeochemical simulations necessitate the discrete integration of mass and energy conservation equations over the spherical surface of the Earth. Conventional regular latitude-longitude grids introduce polar coordinate singularities and drastic cell-area variations, severely penalizing explicit numerical integration time steps through stringent Courant-Friedrichs-Lewy (CFL) stability conditions.

Discrete Global Grid Systems (DGGS), such as the icosahedral hexagonal grid implemented in Uber H3, mitigate these distortions by partitioning the spherical manifold into approximately equal-area Voronoi cells. Nevertheless, modeling cross-boundary physical transport—such as Darcy-Richards porous subsurface groundwater flow, Manning-Strickler diffusive overland runoff, boundary-layer eddy diffusion, and dissolved nutrient advection—requires metric quantities defined along the shared boundary edges between adjacent cells.

Without explicit interface metrics, discrete transport models rely on isotropic point-to-point approximations. These approximations fail to capture anisotropic Voronoi edge lengths, non-orthogonal surface normal orientations, and elevation gradients. Consequently, they introduce unphysical artificial divergence and violate thermodynamic conservation laws.

Sprint 051 introduces the formal contract \texttt{H3CellInterfaceMetrics} within the Web of Life engine (\url{https://github.com/pascalranoroarijaona/WebOfLife}), establishing a verified mathematical foundation for all subsequent finite-volume boundary operators.

\section{Mathematical Formulation}

\subsection{Interface Metric Parameterization}
Let $\mathcal{M}$ denote the Riemannian manifold of the spherical surface, partitioned into discrete H3 cells $\mathcal{C} = \{c_1, c_2, \dots, c_N\}$. For any topologically adjacent pair $(c_i, c_j)$, the interface metric $M_{ij}$ is defined as a 10-tuple:
\begin{equation}
M_{ij} = \langle c_i, c_j, L_{ij}, d_{ij}, \theta_{ij}, \hat{n}_{ij}, A_{ij}^{\text{atm}}, A_{ij}^{\text{sub}}, S_{ij}, \gamma_{ij} \rangle
\end{equation}
where:
\begin{itemize}
    \item $L_{ij} \in \mathbb{R}^+$ is the geodesic length of the shared Voronoi edge in meters.
    \item $d_{ij} \in \mathbb{R}^+$ is the geodesic centroid-to-centroid distance in meters.
    \item $\theta_{ij} \in [0, 2\pi)$ is the geodetic azimuth from $c_i$ to $c_j$ in radians.
    \item $\hat{n}_{ij} \in \mathbb{R}^3$ is the unit normal vector pointing from $c_i$ to $c_j$ in the local tangent plane (East, North, Up).
    \item $A_{ij}^{\text{atm}} \in \mathbb{R}^+$ is the vertical cross-sectional area for atmospheric boundary layer exchange ($\text{m}^2$).
    \item $A_{ij}^{\text{sub}} \in \mathbb{R}^+$ is the vertical cross-sectional area of the subterranean soil and porous media column ($\text{m}^2$).
    \item $S_{ij} \in \mathbb{R}$ is the non-dimensional topographic slope gradient along the interface normal:
    \begin{equation}
    S_{ij} = \frac{z_j - z_i}{d_{ij}}
    \end{equation}
    where $z_i, z_j$ represent cell surface elevations.
    \item $\gamma_{ij} \in \mathbb{R}^+$ is the dimensionless geometric conductance factor:
    \begin{equation}
    \gamma_{ij} = \frac{L_{ij}}{d_{ij}}
    \end{equation}
\end{itemize}

\subsection{Thermodynamic Invariants}

\subsubsection{First Law (Conservation Reciprocity)}
For any conservative physical stock $S \in \{M_{\text{water}}, M_{\text{carbon}}, M_{\text{solute}}, U_{\text{thermal}}\}$, the cross-cell flux $\Phi_{i \to j}(S)$ across $M_{ij}$ must satisfy strict antisymmetry:
\begin{equation}
\Phi_{i \to j}(S) + \Phi_{j \to i}(S) = 0 \quad \forall (c_i, c_j)
\end{equation}
This is guaranteed by the geometric reciprocity of the interface metrics:
\begin{equation}
L_{ij} = L_{ji}, \quad d_{ij} = d_{ji}, \quad \hat{n}_{ij} = -\hat{n}_{ji}, \quad S_{ij} = -S_{ji}
\end{equation}
Hence, integrating over any closed boundary $\partial \Omega$ comprising interface subsets yields zero spurious divergence:
\begin{equation}
\oint_{\partial \Omega} \vec{J} \cdot d\vec{A} = \sum_{c_i \in \Omega} \sum_{c_j \in \mathcal{N}(i) \setminus \Omega} \Phi_{i \to j} = 0
\end{equation}
in the absence of external boundary conditions or localized source terms.

\subsubsection{Second Law (Non-Negative Entropy Production)}
Thermal diffusion across the atmospheric cross-sectional area $A_{ij}^{\text{atm}}$ is modeled as:
\begin{equation}
J_{H, \text{diff}} = - \rho_{\text{air}} c_p K_T \left( \frac{A_{ij}^{\text{atm}}}{d_{ij}} \right) (T_j - T_i)
\end{equation}
where $\rho_{\text{air}} > 0$ is atmospheric air density, $c_p > 0$ is specific heat capacity, and $K_T > 0$ is thermal eddy diffusivity.

The corresponding rate of entropy generation $\dot{S}_{\text{gen}}$ in the isolated two-cell subsystem is:
\begin{equation}
\begin{aligned}
\dot{S}_{\text{gen}} &= \frac{d S_i}{dt} + \frac{d S_j}{dt} = - \frac{J_{H, \text{diff}}}{T_i} + \frac{J_{H, \text{diff}}}{T_j} \\
&= J_{H, \text{diff}} \left( \frac{1}{T_j} - \frac{1}{T_i} \right) \\
&= \rho_{\text{air}} c_p K_T \left( \frac{A_{ij}^{\text{atm}}}{d_{ij}} \right) \frac{(T_j - T_i)^2}{T_i T_j} \ge 0
\end{aligned}
\end{equation}
Because $A_{ij}^{\text{atm}} > 0$, $d_{ij} > 0$, and absolute temperatures $T_i, T_j > 0$, $\dot{S}_{\text{gen}}$ is strictly non-negative, satisfying the Second Law of Thermodynamics.

\section{Software Implementation}
The metric interface is codified in TypeScript within \texttt{src/spatial/h3\_types.ts}:

\begin{lstlisting}
export interface H3CellInterfaceMetrics {
  readonly originIndex: string;
  readonly neighborIndex: string;
  readonly sharedEdgeLengthMeters: number;
  readonly centroidDistanceMeters: number;
  readonly bearingRadians: number;
  readonly normalVector: readonly [number, number, number];
  readonly atmosphericContactAreaM2: number;
  readonly subterraneanContactAreaM2: number;
  readonly topographicSlope: number;
  readonly geometricConductance: number;
}
\end{lstlisting}

Immutable readonly semantics ensure thread-safety and prevent accidental state corruption during concurrent parallel execution on multi-core Node.js runtimes.

\section{Experimental Verification}

The verification test harness is implemented in Node.js and executed using \texttt{npx tsx}:
\begin{lstlisting}
npm install
npx tsx tests/sprint_051.test.ts
\end{lstlisting}

\begin{table}[htbp]
\caption{Invariant Verification Suite Results}
\begin{center}
\begin{tabular}{llcc}
\toprule
\textbf{Invariant ID} & \textbf{Physical Meaning} & \textbf{Tolerance} & \textbf{Result} \\
\midrule
INV-051-A & Edge Length Symmetry & $10^{-12}$ m & PASS \\
INV-051-B & Centroid Distance Symmetry & $10^{-12}$ m & PASS \\
INV-051-C & Normal Vector Inversion & $10^{-9}$ & PASS \\
INV-051-D & Topographic Slope Antisymmetry & $10^{-9}$ & PASS \\
INV-051-E & Conductance Positivity & $> 0$ & PASS \\
INV-051-FLUX & Conservation Antisymmetry & $10^{-9}$ & PASS \\
INV-051-ENT & Non-Negative Entropy & $\ge 0$ & PASS \\
\bottomrule
\end{tabular}
\end{center}
\end{table}

Machine-precision tests confirm that all reciprocal pairs evaluate to zero net divergence across simulated transfer iterations.

\section{Conclusion and Future Work}
The formal definition of \texttt{H3CellInterfaceMetrics} resolves the spatial discretization artifacts associated with point-to-point transfers in Discrete Global Grid Systems. By parameterizing finite-volume boundary operators with exact shared lengths, normal vectors, and stratified contact cross-sections, the Web of Life platform establishes a thermodynamically consistent substrate for global hydrological, atmospheric, and ecological transport. Subsequent sprints will introduce GPU-accelerated computing of adjacency interfaces in \texttt{H3AdjacencyManager} and vectorized Laplacian operators within \texttt{SpatialMonad}.

\begin{thebibliography}{00}
\bibitem{sahr2003} K. Sahr, D. White, and A. J. Kimerling, ``Geodesic discrete global grid systems,'' \textit{Cartography and Geographic Information Science}, vol. 30, no. 2, pp. 121--134, 2003.
\bibitem{uberh3} Uber Technologies, Inc., ``H3: A Hexagonal Hierarchical Spatial Index,'' 2018. [Online]. Available: \url{https://h3geo.org}.
\bibitem{degroot1962} S. R. de Groot and P. Mazur, \textit{Non-Equilibrium Thermodynamics}. North-Holland Publishing Company, Amsterdam, 1962.
\end{thebibliography}

\end{document}